Training of the five red--green--blue (RGB) backbones across the four main experiments is complete, with two additional probes reported below: a frequency-only ResNet-18 baseline for Exp1/Exp2 and a no-retraining score-level temporal aggregation baseline. Unless stated, all metrics are \emph{video-level} (face-crop scores aggregated per video). Accuracy and F1 use a fixed threshold of $0.5$; AUC is threshold-independent and EER is derived from the error-rate curves. Reported 95\% bootstrap intervals apply to AUC.

\subsection{Backbone comparison}
Table~\ref{tab:main_auc} reports video-level AUC for every backbone and experiment, and Table~\ref{tab:main_full} gives the corresponding accuracy, F1 and EER. These values are interpreted as single-seed point estimates; the paired tests below assess held-out-video uncertainty but not training-seed variability. ConvNeXt-Tiny and ResNet-50 yield the two highest observed AUCs in Exp1, Exp2 and Exp4. ConvNeXt-Tiny is highest by point estimate, but its gap over ResNet-50 is not significant under the reported Holm-adjusted procedures in Exp1/Exp2. ViT-Base does not lead under this protocol, and Xception has the lowest in-distribution point estimate. These observations do not establish a general convolutional advantage or a stable architectural ranking. Because the five backbones share a common baseline pipeline, architecture is the deliberately varied factor, but each architecture is not independently optimized. The EfficientNet-B4 ablation further shows that removing progressive unfreezing can materially improve that model.

\begin{table*}[!htbp]
\caption{Video-level AUC (95\% bootstrap CI) per backbone across the four experiments. The best \emph{point estimate} per column is in \textbf{bold}; because several CIs overlap, the bolded value is not always statistically separated from the next model --- see the paired significance tests (Tables~\ref{tab:sig_exp1}--\ref{tab:sig_exp2}). Exp3 shows weak or near-chance transfer, so its column ranking should be interpreted cautiously.}
\centering
\small
\setlength{\tabcolsep}{3pt}
\resizebox{\textwidth}{!}{%
\begin{tabular}{l cccc}
\toprule
\textbf{Backbone} & \textbf{Exp1: FF++$\to$FF++} & \textbf{Exp2: FF++$\to$Celeb-DF} & \textbf{Exp3: Celeb-DF$\to$FF++} & \textbf{Exp4: Mixed$\to$DFD} \\
\midrule
EfficientNet-B4 & 0.900 (0.85--0.94) & 0.713 (0.63--0.80) & 0.506 (0.47--0.54) & 0.788 (0.75--0.82) \\
Xception & 0.736 (0.65--0.81) & 0.573 (0.47--0.67) & 0.543 (0.51--0.57) & 0.735 (0.70--0.77) \\
ResNet-50 & 0.924 (0.88--0.97) & 0.799 (0.72--0.87) & \textbf{0.589 (0.56--0.62)} & 0.812 (0.78--0.84) \\
ConvNeXt-Tiny & \textbf{0.957 (0.92--0.99)} & \textbf{0.806 (0.72--0.88)} & 0.576 (0.54--0.61) & \textbf{0.915 (0.89--0.93)} \\
ViT-Base & 0.867 (0.80--0.93) & 0.668 (0.57--0.76) & 0.588 (0.56--0.62) & 0.713 (0.67--0.75) \\
\botrule
\end{tabular}%
}
\label{tab:main_auc}
\end{table*}

\begin{table*}[!htbp]
\caption{Video-level Accuracy, F1 and EER per backbone, grouped by experiment (lower EER is better).}
\centering
\small
\setlength{\tabcolsep}{3pt}
\begin{tabular}{l l ccc}
\toprule
\textbf{Experiment} & \textbf{Backbone} & \textbf{Acc} & \textbf{F1} & \textbf{EER} \\
\midrule
Exp1: FF++$\to$FF++ & EfficientNet-B4 & 0.799 & 0.741 & 0.194 \\
 & Xception & 0.671 & 0.680 & 0.322 \\
 & ResNet-50 & 0.859 & 0.835 & 0.174 \\
 & ConvNeXt-Tiny & 0.886 & 0.866 & 0.087 \\
 & ViT-Base & 0.765 & 0.759 & 0.202 \\
\midrule
Exp2: FF++$\to$Celeb-DF & EfficientNet-B4 & 0.659 & 0.633 & 0.324 \\
 & Xception & 0.535 & 0.552 & 0.442 \\
 & ResNet-50 & 0.744 & 0.680 & 0.227 \\
 & ConvNeXt-Tiny & 0.767 & 0.727 & 0.240 \\
 & ViT-Base & 0.659 & 0.569 & 0.372 \\
\midrule
Exp3: Celeb-DF$\to$FF++ & EfficientNet-B4 & 0.498 & 0.287 & 0.490 \\
 & Xception & 0.484 & 0.232 & 0.455 \\
 & ResNet-50 & 0.516 & 0.152 & 0.445 \\
 & ConvNeXt-Tiny & 0.523 & 0.094 & 0.449 \\
 & ViT-Base & 0.528 & 0.167 & 0.444 \\
\midrule
Exp4: Mixed$\to$DFD & EfficientNet-B4 & 0.688 & 0.598 & 0.271 \\
 & Xception & 0.614 & 0.686 & 0.331 \\
 & ResNet-50 & 0.694 & 0.582 & 0.243 \\
 & ConvNeXt-Tiny & 0.710 & 0.589 & 0.159 \\
 & ViT-Base & 0.650 & 0.602 & 0.338 \\
\botrule
\end{tabular}
\label{tab:main_full}
\end{table*}

\subsection{Statistical significance of backbone differences}
Because all results are single-seed, we test whether the point-estimate gaps in Table~\ref{tab:main_auc} are statistically supported before drawing architectural conclusions. For Exp1 and Exp2 and each pair of backbones we compute, on the common set of test videos, (i)~a \emph{paired bootstrap} on the AUC difference ($2000$ resamples over videos, with a $95\%$ CI and an empirical two-sided bootstrap-tail value) and (ii)~an \emph{exact McNemar test} on the binary decisions at the fixed $0.5$ threshold. The pairwise outcomes for the two experiments with clear signal are reported in Table~\ref{tab:sig_exp1} (in-distribution Exp1) and Table~\ref{tab:sig_exp2} (cross-dataset FF++$\to$Celeb-DF). \textbf{Scope of these tests:} the paired bootstrap resamples \emph{test videos}, whereas exact McNemar tests use discordant binary decisions; they capture test-set sampling uncertainty \emph{at a single training seed only}; they do \emph{not} account for training/initialisation (seed) variance, which would require retraining and is out of scope here. To reduce false discoveries across the ten pairwise backbone comparisons in each reported experiment, raw bootstrap and McNemar p-values are accompanied by Holm-Bonferroni adjusted q-values, corrected separately for each experiment and test family. The tests are computed on a regenerated, seeded five-view TTA per-video score set whose individual AUCs differ from the original Table~\ref{tab:main_auc} evaluations by at most $0.003210$ AUC. Stochastic brightness TTA is a possible, but unconfirmed, explanation. Tables~S3--S7 use this regenerated-score evaluation, not the original headline evaluation; pairwise $\Delta$AUC values therefore need not equal differences between the headline entries.

The original generating script was not recovered. A reconstructed numerical procedure reproduces the stored bootstrap intervals and tail values: align and sort by video identifier, reset NumPy's \texttt{default\_rng(42)} for each model pair, and draw 2,000 paired video samples with replacement. The CI uses the 2.5th and 97.5th percentiles of $\Delta_b$, the resampled AUC difference. The reported $p_{\mathrm{boot}}$ is twice the smaller fraction of strictly negative and strictly positive differences, capped at one; exact zeros count in neither tail. Stored zeros, displayed as $<0.001$, mean no draws entered the relevant strict tail, not zero underlying probability. Numerical reproduction does not establish calibration of this tail calculation as a hypothesis test; bootstrap-based significance labels are conditional on this convention.

Under this convention, only the larger gaps survive correction. In the in-distribution setting, ResNet-50 and ConvNeXt-Tiny are each stronger than Xception on the AUC bootstrap and McNemar tests after Holm adjustment; EfficientNet-B4 remains separated from Xception on the AUC bootstrap, while its McNemar result is only nominal. ViT-Base is also stronger than Xception on AUC ($\Delta$AUC $=0.131$, $q_{\mathrm{boot}}\!=\!0.007$), but not on the fixed-threshold McNemar test ($q_{\mathrm{McN}}\!=\!0.304$). ConvNeXt-Tiny's gap over EfficientNet-B4 is nominal before correction ($p_{\mathrm{boot}}\!=\!0.020$, $p_{\mathrm{McN}}\!=\!0.007$) but borderline after correction ($q_{\mathrm{boot}}\!=\!0.100$, $q_{\mathrm{McN}}\!=\!0.0504$), so it should not be read as a decisive adjusted result. The comparison that matters most, ConvNeXt-Tiny versus ResNet-50, is not significant after adjustment in Exp1 ($q_{\mathrm{boot}}\!=\!0.147$, $q_{\mathrm{McN}}\!=\!0.688$) and is not statistically significant in Exp2 ($\Delta$AUC $=0.007$, $q_{\mathrm{boot}}\!=\!0.875$). I report ConvNeXt-Tiny as the strongest observed point estimate, not as a uniquely superior architecture. Pairwise tests for Exp3 and Exp4 are not reported. Full $\Delta$AUC/CI/$p/q$ values are listed in the significance tables.

\begin{table*}[!htbp]
\caption{Pairwise backbone comparison from regenerated scores on Exp1: FF++$\to$FF++ (149 videos). $\Delta$AUC $=$ AUC(A)$-$AUC(B) with paired-bootstrap 95\% CI; $p_{\mathrm{boot}}$ and $p_{\mathrm{McN}}$ are raw empirical bootstrap-tail values and exact McNemar p-values, while $q_{\mathrm{boot}}$ and $q_{\mathrm{McN}}$ are Holm-Bonferroni adjusted within each experiment/test family. Bold marks Holm-adjusted $q<0.05$.}
\centering
\small
\setlength{\tabcolsep}{3pt}
\begin{tabular}{l l r r r r r}
\toprule
\textbf{A} & \textbf{B} & \textbf{$\Delta$AUC (95\% CI)} & \textbf{$p_{\mathrm{boot}}$} & \textbf{$q_{\mathrm{boot}}$} & \textbf{$p_{\mathrm{McN}}$} & \textbf{$q_{\mathrm{McN}}$} \\
\midrule
EfficientNet-B4 & Xception & \textbf{+0.163 [+0.088, +0.249]} & $<0.001$ & \textbf{$<0.001$} & 0.014 & 0.075 \\
EfficientNet-B4 & ResNet-50 & -0.022 [-0.068, +0.018] & 0.262 & 0.444 & 0.078 & 0.304 \\
EfficientNet-B4 & ConvNeXt-Tiny & -0.059 [-0.110, -0.010] & 0.020 & 0.100 & 0.007 & 0.0504 \\
EfficientNet-B4 & ViT-Base & +0.032 [-0.018, +0.087] & 0.222 & 0.444 & 0.511 & 0.688 \\
Xception & ResNet-50 & \textbf{-0.186 [-0.273, -0.111]} & $<0.001$ & \textbf{$<0.001$} & $<0.001$ & \textbf{$<0.001$} \\
Xception & ConvNeXt-Tiny & \textbf{-0.222 [-0.311, -0.151]} & $<0.001$ & \textbf{$<0.001$} & $<0.001$ & \textbf{$<0.001$} \\
Xception & ViT-Base & \textbf{-0.131 [-0.223, -0.046]} & 0.001 & \textbf{0.007} & 0.076 & 0.304 \\
ResNet-50 & ConvNeXt-Tiny & -0.037 [-0.076, -0.000] & 0.049 & 0.147 & 0.344 & 0.688 \\
ResNet-50 & ViT-Base & +0.055 [+0.005, +0.108] & 0.031 & 0.124 & 0.013 & 0.075 \\
ConvNeXt-Tiny & ViT-Base & \textbf{+0.091 [+0.035, +0.154]} & 0.001 & \textbf{0.007} & 0.001 & \textbf{0.011} \\
\botrule
\end{tabular}
\label{tab:sig_exp1}
\end{table*}

\vspace{0.8em}

\begin{table*}[!htbp]
\caption{Pairwise backbone comparison from regenerated scores on Exp2: FF++$\to$Celeb-DF (129 videos). $\Delta$AUC $=$ AUC(A)$-$AUC(B) with paired-bootstrap 95\% CI; $p_{\mathrm{boot}}$ and $p_{\mathrm{McN}}$ are raw empirical bootstrap-tail values and exact McNemar p-values, while $q_{\mathrm{boot}}$ and $q_{\mathrm{McN}}$ are Holm-Bonferroni adjusted within each experiment/test family. Bold marks Holm-adjusted $q<0.05$.}
\centering
\small
\setlength{\tabcolsep}{3pt}
\begin{tabular}{l l r r r r r}
\toprule
\textbf{A} & \textbf{B} & \textbf{$\Delta$AUC (95\% CI)} & \textbf{$p_{\mathrm{boot}}$} & \textbf{$q_{\mathrm{boot}}$} & \textbf{$p_{\mathrm{McN}}$} & \textbf{$q_{\mathrm{McN}}$} \\
\midrule
EfficientNet-B4 & Xception & +0.144 [+0.023, +0.266] & 0.019 & 0.114 & 0.018 & 0.140 \\
EfficientNet-B4 & ResNet-50 & -0.084 [-0.160, -0.001] & 0.044 & 0.176 & 0.099 & 0.358 \\
EfficientNet-B4 & ConvNeXt-Tiny & -0.092 [-0.175, -0.008] & 0.026 & 0.130 & 0.035 & 0.210 \\
EfficientNet-B4 & ViT-Base & +0.047 [-0.048, +0.145] & 0.329 & 0.658 & 1.000 & 1.000 \\
Xception & ResNet-50 & \textbf{-0.228 [-0.353, -0.099]} & $<0.001$ & \textbf{$<0.001$} & $<0.001$ & \textbf{0.004} \\
Xception & ConvNeXt-Tiny & \textbf{-0.236 [-0.369, -0.102]} & 0.001 & \textbf{0.009} & $<0.001$ & \textbf{0.002} \\
Xception & ViT-Base & -0.097 [-0.232, +0.036] & 0.144 & 0.432 & 0.046 & 0.232 \\
ResNet-50 & ConvNeXt-Tiny & -0.007 [-0.076, +0.063] & 0.875 & 0.875 & 0.690 & 1.000 \\
ResNet-50 & ViT-Base & \textbf{+0.131 [+0.048, +0.216]} & 0.002 & \textbf{0.016} & 0.090 & 0.358 \\
ConvNeXt-Tiny & ViT-Base & \textbf{+0.138 [+0.052, +0.224]} & 0.003 & \textbf{0.021} & 0.020 & 0.140 \\
\botrule
\end{tabular}
\label{tab:sig_exp2}
\end{table*}


\subsection{Cross-dataset generalization}
Every backbone degrades from the in-distribution setting (Exp1) to the cross-dataset settings, so the in-domain numbers overstate deployment performance. The FF++$\to$Celeb-DF (Exp2) and Mixed$\to$DFD (Exp4) transfers retain useful signal (AUC $0.57$--$0.92$), but \textbf{the reverse Celeb-DF$\to$FF++ transfer (Exp3) collapses to near chance ($\sim$0.51--0.59) for \emph{every} backbone}. Within this protocol, the asymmetry is consistent with source-domain scale and diversity contributing to the cross-dataset result: models trained on the larger, more varied FF++ transfer better than models trained on Celeb-DF and evaluated on FF++'s manipulation families. However, the present experiments do not isolate dataset scale and diversity from other differences between the datasets, so this interpretation should not be treated as a causal finding. The vector summary in the main paper visualizes the same AUC pattern.

\subsection{Per-manipulation analysis (Exp1)}
We report the manipulation-family breakdown for the in-distribution FF++ test set in Exp1. For each backbone and each family we report the video-level AUC of that family's fakes against the shared pool of real videos (Table~\ref{tab:per_manip}); a lower value marks a manipulation that the trained model finds harder to detect. NeuralTextures has the lowest observed family AUC for four of the five backbones ($0.66$--$0.87$), while Xception's lowest value occurs for Face2Face ($0.60$). Under the fixed seed and protocol, the family-difficulty ordering is broadly similar across backbones. ConvNeXt-Tiny has the highest observed point estimate for every family ($0.87$--$0.99$), but these family-level rankings are descriptive and were not tested for stability across training seeds.

\begin{table*}[!htbp]
\caption{Per-manipulation detection from regenerated scores on Exp1 (FF++ in-distribution): video-level AUC of each manipulation family's fakes against the shared real videos. Lower values mark manipulation types a backbone finds harder.}
\centering
\small
\setlength{\tabcolsep}{3pt}
\begin{tabular}{lccccc}
\toprule
\textbf{Backbone} & \textbf{Deepfakes} & \textbf{Face2Face} & \textbf{FaceShifter} & \textbf{FaceSwap} & \textbf{NeuralTextures} \\
\midrule
EfficientNet-B4 & 0.940 & 0.887 & 0.912 & 0.915 & 0.810 \\
Xception & 0.860 & 0.599 & 0.652 & 0.850 & 0.697 \\
ResNet-50 & 0.953 & 0.940 & 0.851 & 0.974 & 0.819 \\
ConvNeXt-Tiny & 0.973 & 0.983 & 0.924 & 0.985 & 0.874 \\
ViT-Base & 0.934 & 0.862 & 0.861 & 0.930 & 0.656 \\
\botrule
\end{tabular}
\label{tab:per_manip}
\end{table*}


\subsection{Error analysis (Confusion Matrices)}
Table~\ref{tab:confusion_exp1} shows the Exp1 confusion counts at a threshold of $0.5$ ($80$ real and $69$ fake videos). ConvNeXt-Tiny, EfficientNet-B4, and ResNet-50 have the three lowest false-positive counts, with FP $=3$, $4$, and $5$, respectively. EfficientNet-B4 misses 26 fake videos. ViT-Base has more false positives (FP $=21$), while Xception has the largest number, with 32 of the 80 real videos classified as fake. Xception has 17 false negatives.

Looking at these counts gives a somewhat different picture from AUC alone. The relevant error can depend on how the detector is used. In a moderation queue, missing a fake may be more concerning, whereas in forensic triage, a high number of false positives could mean that more real videos have to be checked manually.

\begin{table*}[!htbp]
\caption{Video-level confusion counts from regenerated scores for Exp1 on FF++ (80 real and 69 fake videos; threshold $0.5$). TP and TN denote correctly classified fake and real videos, respectively. FP denotes a real video classified as fake, and FN denotes a fake video classified as real.}
\centering
\small
\setlength{\tabcolsep}{3pt}
\begin{tabular}{l cccc}
\toprule
\textbf{Backbone} & \textbf{TP} & \textbf{TN} & \textbf{FP} & \textbf{FN} \\
\midrule
EfficientNet-B4 & 43 & 76 & 4 & 26 \\
Xception & 52 & 48 & 32 & 17 \\
ResNet-50 & 53 & 75 & 5 & 16 \\
ConvNeXt-Tiny & 55 & 77 & 3 & 14 \\
ViT-Base & 55 & 59 & 21 & 14 \\
\botrule
\end{tabular}
\label{tab:confusion_exp1}
\end{table*}


\subsection{Deeper Failure-Case Analysis}

The saved video predictions allow the errors to be counted directly for Exp1 and Exp2 (Table~\ref{tab:error_deep_dive}). In Exp1, ConvNeXt-Tiny makes 17 errors out of 149 videos, compared with 21 for ResNet-50. Xception has a much higher error count, with 49 misclassified videos, most of them false positives.

Error counts increase for every backbone in Exp2. ConvNeXt-Tiny makes 30 errors out of 129 videos, while ResNet-50 makes 33. The two remain close in terms of error count, consistent with the paired test results. EfficientNet-B4 and ViT-Base both have error rates close to one-third. Xception again produces a relatively large number of false positives on real videos.


\begin{table*}[!htbp]
\caption{Failure-case summary from regenerated scores for Exp1 and Exp2 at the fixed 0.5 video-level threshold. FP and FN denote false positives and false negatives, respectively.}
\centering
\small
\setlength{\tabcolsep}{3pt}
\begin{tabular}{l l c c c c}
\toprule
\textbf{Experiment} & \textbf{Backbone} & \textbf{FP} & \textbf{FN} & \textbf{Total errors} & \textbf{Error rate} \\
\midrule
Exp1: FF++$\to$FF++ & EfficientNet-B4 & 4 & 26 & 30/149 & 0.201 \\
 & Xception & 32 & 17 & 49/149 & 0.329 \\
 & ResNet-50 & 5 & 16 & 21/149 & 0.141 \\
 & ConvNeXt-Tiny & \textbf{3} & \textbf{14} & \textbf{17/149} & \textbf{0.114} \\
 & ViT-Base & 21 & 14 & 35/149 & 0.235 \\
\midrule
Exp2: FF++$\to$Celeb-DF & EfficientNet-B4 & 28 & 15 & 43/129 & 0.333 \\
 & Xception & 44 & 17 & 61/129 & 0.473 \\
 & ResNet-50 & \textbf{15} & 18 & 33/129 & 0.256 \\
 & ConvNeXt-Tiny & 16 & \textbf{14} & \textbf{30/129} & \textbf{0.233} \\
 & ViT-Base & 19 & 25 & 44/129 & 0.341 \\
\botrule
\end{tabular}
\label{tab:error_deep_dive}
\end{table*}

The manipulation breakdown shows where many Exp1 errors occur. NeuralTextures is particularly difficult for the three strongest backbones: ConvNeXt-Tiny misses 5/9 videos, ResNet-50 misses 6/9, and EfficientNet-B4 misses 7/9. Real videos show a different pattern. ConvNeXt-Tiny and ResNet-50 produce only 3/80 and 5/80 false positives, respectively, while ViT-Base and Xception produce 21/80 and 32/80.

\subsection{Ablations (EfficientNet-B4)}
Tables~\ref{tab:ablation} and~\ref{tab:agg} ablate the training and inference choices on the EfficientNet-B4 baseline (video-level AUC). Label smoothing has negligible effect; TTA improves AUC in Exp1/Exp2 but slightly reduces it in Exp3/Exp4. Training augmentation helps, with the largest observed gain in Exp4 ($0.79$ vs.\ $0.62$ without it). For video aggregation, mean and median are best and clearly beat max/top-$k$, justifying mean as the default.

The most consequential finding is on \emph{progressive unfreezing}. Removing it (i.e.\ fine-tuning all stages from the start) \textbf{improves EfficientNet-B4 in every experiment} --- Exp1 $0.900\!\to\!0.927$, Exp2 $0.713\!\to\!0.812$, Exp3 $0.506\!\to\!0.546$, Exp4 $0.788\!\to\!0.813$ --- an observed configuration-level improvement that does not isolate freezing from the accompanying optimizer and learning-rate changes. I identify the \textbf{no-progressive-unfreeze recipe as the recommended EfficientNet-B4 configuration}, and report it explicitly in the ablation. I deliberately retain the common recipe in the headline backbone comparison (Table~\ref{tab:main_auc}) because only EfficientNet-B4 was ablated; re-running the better recipe for the other four backbones would require additional training, which is outside the scope of this no-retraining study and is flagged in future work. The cross-configuration ranking changes: EfficientNet-B4 overtakes ResNet-50 in Exp1/Exp4 and both ResNet-50 and ConvNeXt-Tiny in Exp2, although it remains below ConvNeXt-Tiny in Exp1. These single-seed comparisons do not establish architectural superiority.

\begin{table*}[!htbp]
\caption{Training ablations (EfficientNet-B4, video-level AUC). ``$-$X'' removes component X from the full recipe; experiment abbreviations follow Table~\ref{tab:main_auc}.}
\centering
\small
\setlength{\tabcolsep}{3pt}
{\renewcommand{\arraystretch}{1.08}
\begin{tabular}{l cccc}
\toprule
\textbf{Setting} & \textbf{Exp1} & \textbf{Exp2} & \textbf{Exp3} & \textbf{Exp4} \\
\midrule
Full (baseline) & 0.900 & 0.713 & 0.506 & 0.788 \\
-- label smoothing & 0.898 & 0.708 & 0.505 & 0.790 \\
-- progressive unfreeze & 0.927 & 0.812 & 0.546 & 0.813 \\
-- train augmentation & 0.888 & 0.704 & 0.518 & 0.619 \\
-- TTA & 0.890 & 0.690 & 0.508 & 0.794 \\
\botrule
\end{tabular}}
\label{tab:ablation}
\end{table*}

\begin{table*}[!htbp]
\caption{Video-score aggregation ablation (EfficientNet-B4, video-level AUC); experiment abbreviations follow Table~\ref{tab:main_auc}.}
\centering
\small
\setlength{\tabcolsep}{3pt}
{\renewcommand{\arraystretch}{1.08}
\begin{tabular}{l cccc}
\toprule
\textbf{Setting} & \textbf{Exp1} & \textbf{Exp2} & \textbf{Exp3} & \textbf{Exp4} \\
\midrule
Mean (baseline) & 0.900 & 0.713 & 0.506 & 0.788 \\
Median & 0.899 & 0.700 & 0.506 & 0.789 \\
Max & 0.861 & 0.642 & 0.492 & 0.759 \\
Top-$k$ & 0.883 & 0.677 & 0.497 & 0.775 \\
\botrule
\end{tabular}}
\label{tab:agg}
\end{table*}

\subsection{Frequency-only and temporal-score probes}
Two additional low-cost baselines test whether simple alternatives to RGB mean aggregation reduce the transfer problem. The first is a frequency-only ResNet-18 trained on DCT/log-magnitude face crops for Exp1 and Exp2. The second is a score-level temporal baseline that trains a lightweight logistic-regression classifier over per-video score statistics, using existing crop-score CSVs without retraining the image models. Exp2 fits on Exp1 scores and tests on scores from separately trained Exp2 detectors of the same architecture; it is exploratory cross-checkpoint transfer, not a fixed-detector aggregation comparison.

\begin{table*}[!htbp]
\caption{Lightweight frequency-only ResNet-18 baseline. Inputs are discrete cosine transform (DCT)/log-magnitude face crops, not RGB images. Values are video-level metrics with AUC bootstrap CI.}
\centering
\small
\setlength{\tabcolsep}{3pt}
\begin{tabular}{l c c c c}
\toprule
\textbf{Experiment} & \textbf{AUC (95\% CI)} & \textbf{Acc} & \textbf{F1} & \textbf{EER} \\
\midrule
Exp1: FF++$\to$FF++ & 0.653 (0.561--0.733) & 0.617 & 0.544 & 0.409 \\
Exp2: FF++$\to$Celeb-DF & 0.752 (0.657--0.835) & 0.519 & 0.622 & 0.259 \\
\botrule
\end{tabular}
\label{tab:freq_resnet18}
\end{table*}

The frequency result (Table~\ref{tab:freq_resnet18}) is mixed rather than a clear improvement. In Exp1, its AUC is well below that of the strongest RGB backbones in this study ($0.653$ versus $0.957$ for ConvNeXt-Tiny and $0.924$ for ResNet-50). In Exp2, it reaches an AUC of $0.752$, higher than EfficientNet-B4 and ViT-Base but lower than ConvNeXt-Tiny and ResNet-50. The low accuracy at the fixed $0.5$ threshold indicates that the AUC in Exp2 does not translate into similarly strong classification at that threshold. Spectral input appears to carry some transfer signal, but this simple ResNet-18 probe does not match the stronger RGB backbones.

\begin{table*}[!htbp]
\caption{Score-level temporal aggregation probe. ``Mean'' is the matched video mean score; ``Temporal'' is logistic regression over per-video score statistics. Exp1 fits on 104 videos and evaluates on 45 (an internal 70/30 split); Exp2 fits on Exp1 FF++ scores and tests on separately trained Exp2 detector scores from Celeb-DF (exploratory cross-checkpoint transfer).}
\centering
\small
\setlength{\tabcolsep}{3pt}
\begin{tabular}{l l c c c}
\toprule
\textbf{Experiment} & \textbf{Backbone} & \textbf{Mean AUC} & \textbf{Temporal AUC} & \textbf{$\Delta$AUC} \\
\midrule
Exp1: FF++$\to$FF++ & EfficientNet-B4 & 0.823 & 0.704 & $-$0.119 \\
 & Xception & 0.754 & 0.625 & $-$0.129 \\
 & ResNet-50 & 0.873 & 0.847 & $-$0.026 \\
 & ConvNeXt-Tiny & 0.968 & 0.915 & $-$0.054 \\
 & ViT-Base & 0.841 & 0.794 & $-$0.048 \\
\midrule
Exp2: FF++$\to$Celeb-DF & EfficientNet-B4 & 0.714 & 0.725 & $+$0.011 \\
 & Xception & 0.570 & 0.557 & $-$0.013 \\
 & ResNet-50 & 0.798 & 0.786 & $-$0.012 \\
 & ConvNeXt-Tiny & 0.806 & 0.747 & $-$0.059 \\
 & ViT-Base & 0.667 & 0.642 & $-$0.026 \\
\botrule
\end{tabular}
\label{tab:temporal_score_probe}
\end{table*}

\begin{table*}[!htbp]
\caption{Robustness of EfficientNet-B4: clean video-level AUC and minimum observed AUC across tested severities of each corruption (drop in parentheses).}
\centering
\small
\setlength{\tabcolsep}{3pt}
\begin{tabular}{l c cccc}
\toprule
\textbf{Experiment} & \textbf{Clean} & \textbf{JPEG} & \textbf{Blur} & \textbf{Noise} & \textbf{Downscale} \\
\midrule
Exp1: FF++$\to$FF++ & 0.890 & 0.734 ($-$0.156) & 0.740 ($-$0.150) & 0.476 ($-$0.414) & 0.780 ($-$0.109) \\
Exp2: FF++$\to$Celeb-DF & 0.690 & 0.640 ($-$0.050) & 0.560 ($-$0.130) & 0.465 ($-$0.225) & 0.588 ($-$0.101) \\
Exp3: Celeb-DF$\to$FF++ & 0.508 & 0.502 ($-$0.006) & 0.495 ($-$0.013) & 0.459 ($-$0.049) & 0.491 ($-$0.017) \\
Exp4: Mixed$\to$DFD & 0.793 & 0.692 ($-$0.101) & 0.589 ($-$0.205) & 0.489 ($-$0.304) & 0.648 ($-$0.145) \\
\botrule
\end{tabular}
\label{tab:robust}
\end{table*}

The temporal-score probe (Table~\ref{tab:temporal_score_probe}) is also mostly negative. It gives a small Exp2 AUC gain for EfficientNet-B4 ($0.714\to0.725$), but it reduces AUC for the other four backbones in Exp2 and for all five backbones in the Exp1 split. The internal Exp1 split shows no AUC improvement, but Exp2 confounds dataset shift with a detector-checkpoint change. These results do not rule out mean aggregation as a limitation. A fixed-detector test would require source-calibration and target-test scores from the same checkpoint.

\subsection{Robustness to image degradations}
Table~\ref{tab:robust} reports the minimum observed EfficientNet-B4 video-level AUC across the tested severities of each corruption. To keep the comparison clean, both the clean baseline and the corrupted evaluations use \emph{single-view} (no-TTA) inference; each clean baseline comes from its corresponding robustness evaluation, not the separate no-TTA ablation or the TTA headline evaluation. \textbf{Additive noise is by far the most damaging corruption in every experiment} (maximum observed AUC drops of $0.41$/$0.23$/$0.30$ on Exp1/Exp2/Exp4), collapsing the detector toward chance. The \emph{milder} corruptions, however, do \emph{not} order consistently across experiments: JPEG compression is the mildest on the two cross-dataset transfers (Exp2 drop $0.05$, Exp4 drop $0.10$), but on the in-distribution Exp1 \emph{downscaling} is the mildest (drop $0.11$) and JPEG is in fact the second-most damaging after noise (drop $0.16$). Even the DFD detector (clean AUC $0.79$) falls to $\sim$0.49 under heavy noise, showing a measured robustness failure instead of a broad resilience claim. (Exp3 is omitted from interpretation because its clean baseline is already near chance.)
